\subsection{注记与进一步结果}
\label{sec:ch5-notes-further-results}

\subsubsection{参考文献}
\label{subsec:ch5-references}

Calderón--Zygmund 分解及其在证明奇异积分弱型 $(1,1)$ 不等式中的应用, 首见两位作者的经典论文 \MBUnresolvedReference{citation}{[On the existence of certain singular integrals, Acta Math. 88 (1952), 85--139]}. 该文使用梯度条件 \eqref{eq:ch5-gradient-condition}; L. Hörmander 首先指出 \eqref{eq:ch5-hormander-condition} 已经充分, 见 \MBUnresolvedReference{citation}{[Estimates for translation invariant operators in $L^p$-spaces, Acta Math. 104 (1960), 93--139]}. 另见 Stein \MBUnresolvedReference{citation}{[15]}.

命题~\ref{prop:ch5-truncated-kernel-fourier-bound} 归功于 A. Benedek, A. P. Calderón 与 R. Panzone 的 \MBUnresolvedReference{citation}{[Convolution operators with Banach-space valued functions, Proc. Nat. Acad. Sci. U.S.A. 48 (1962), 356--365]}. 分布 $\operatorname{p.v.}|x|^{-n-it}$ 及其推广最早由 B. Muckenhoupt 研究. 关于广义 Calderón--Zygmund 算子, 见 Coifman 与 Meyer \MBUnresolvedReference{citation}{[2]}, Journé \MBUnresolvedReference{citation}{[8]} 以及 Stein \MBUnresolvedReference{citation}{[17]}. Cauchy 积分的相关历史见 N. I. Muskhelishvili 的 \MBUnresolvedReference{citation}{[Singular Integral Equations, P. Noordhoff, Groningen, 1953]}. 相应交换子最早由 A. P. Calderón 研究, 并将在\MBUnresolvedReference{chapter}{第 9 章}再次讨论.

关于向量值奇异积分, 见 García-Cuerva 与 Rubio de Francia \MBUnresolvedReference{citation}{[6, Chapter 5]}, J. L. Rubio de Francia, F. J. Ruiz 与 J. L. Torrea 的论文, 或 J. L. Torrea 的专著. C. Fefferman 的综述 \MBUnresolvedReference{citation}{[Recent progress in classical Fourier analysis, Proceedings of the I.C.M., 1974]} 简要讨论了奇异积分及其与分析中许多其他问题的联系.

\subsubsection{更一般的 Dini 型条件}
\label{subsec:ch5-general-dini}

\hypertarget{chapter-05-page-120}{}
对 $\rho\in O(n)$, 令 $\|\rho\|=\sup\{|u-\rho u|:u\in S^{n-1}\}$, 并定义
\[
 \omega_1(t)=\sup_{\|\rho\|\le t}
 \int_{S^{n-1}}|\Omega(\rho u)-\Omega(u)|\,d\sigma(u).
\]
则更一般的 Dini 型条件
\[
 \int_0^1\omega_1(t)\frac{dt}{t}<\infty
\]
蕴含 $K(x)=\Omega(x')/|x|^n$ 的 Hörmander 条件 \eqref{eq:ch5-hormander-condition}. 这一条件还蕴含 $\Omega\in L\log^+L$. 结果归功于 A. P. Calderón, M. Weiss 与 A. Zygmund; 关于对应条件的必要性, 见 S. Hofmann 的工作.

\subsubsection{非各向同性伸缩}
\label{subsec:ch5-anisotropic-dilations}

Euclidean 范数与伸缩 $\delta_t(x)=tx$ 相联系. 更一般地, 给定 $a_j>0$, 可令
\[
 \delta_t(x)=(t^{a_1}x_1,\ldots,t^{a_n}x_n).
\]
存在相应的拟范数 $\|x\|$, 满足 $\|\delta_t(x)\|=t\|x\|$. 它并非真正的范数, 因为三角不等式只在常数倍意义下成立: $\|x+y\|\le L(\|x\|+\|y\|)$. 在此拟范数下, 把 \eqref{eq:ch5-generalized-hormander-y} 一类条件中的 Euclidean 距离替换为 $\|\cdot\|$, 可以得到定理~\ref{thm:ch5-generalized-calderon-zygmund-bounds} 的类似结论. 这类奇异积分主要出现于与抛物型算子相联系的问题.

\hypertarget{chapter-05-page-121}{}
\subsubsection{齐型空间}
\label{subsec:ch5-spaces-homogeneous-type}

奇异积分还可以在很一般的空间上研究. 设 $X$ 带有拟度量 $\rho$, 即 $\rho:X\times X\to[0,\infty)$ 满足: $\rho(x,y)=0$ 当且仅当 $x=y$; $\rho(x,y)=\rho(y,x)$; 并且 $\rho(x,z)\le L(\rho(x,y)+\rho(y,z))$. 若球 $B(x,r)$ 构成 $X$ 的拓扑基, 并且 $X$ 上有 Borel 测度 $\mu$ 满足倍增条件 $0<\mu(B(x,r))\le A\mu(B(x,r/2))<\infty$, 则称 $X$ 为齐型空间.

若 $K(x,y)$ 定义在 $X\times X\setminus\Delta$ 上, 且
\[
 Tf(x)=\int_XK(x,y)f(y)\,d\mu(y),
\]
若 $T$ 在某个 $L^q$, $q>1$, 上有界, 并满足
\begin{equation}
\label{eq:ch5-homogeneous-space-hormander}
 \int_{\rho(x,y)>2L\rho(y,z)}|K(x,y)-K(x,z)|\,d\mu(x)\le C,
\end{equation}
则 $T$ 是强型 $(p,p)$ 的, $1<p\le q$, 并且是弱型 $(1,1)$ 的. 近年来这一理论被用于 Cauchy 积分, 边界正则性问题和齐流形上的几何积分.

\hypertarget{chapter-05-page-122}{}
上述结果在 $\rho$ 只是度量并满足体积增长 $\mu(B(x,r))\le r^\nu$ 时也成立, 此时应将 \eqref{eq:ch5-standard-kernel-y-regularity} 换成相应的体积归一化估计.

\subsubsection{常数的大小}
\label{subsec:ch5-size-of-constants}

在定理~\ref{thm:ch5-calderon-zygmund} 的证明中, 若在高度 $A^{-1}\lambda$ 作 Calderón--Zygmund 分解, 则弱型常数成为 $C_n(A+B)$, 插值所得的强型常数在 $p\to1$ 或 $p\to\infty$ 时由
\[
 \frac{C_n'(A+B)p^2}{p-1}
\]
控制. 一般说来, 这是最佳渐近阶. 例如取满足梯度条件的核 $K(x)=c|x|^{-n}$ 型模型并考察同心球的示性函数, 可得 $\|T\|_p\ge c(p-1)^{-1}$ 当 $p\downarrow1$; 由对偶性得到 $p\to\infty$ 的相应下界. 对 Hilbert 变换可参见第~\ref{subsec:ch3-size-of-constants}~节的精确常数.

\subsubsection{关于截断积分的进一步结果}
\label{subsec:ch5-more-truncated-integrals}

设 $T$ 是 Calderón--Zygmund 奇异积分, 其截断为 $T_\varepsilon$. 若 $1<p<\infty$, 则由定理~\ref{thm:ch5-maximal-cz-bounds} 和控制收敛定理, $T_\varepsilon f\to Tf$ 在 $L^p$ 中成立. 当 $p=1$ 时同样的论证失效, 但若 $f,Tf\in L^1$, 仍可证明 $T_\varepsilon f\to Tf$ 在 $L^1$ 中成立. 这一结果归功于 A. Calderón 与 O. Capri.

\hypertarget{chapter-05-page-123}{}
\subsubsection{向量值奇异积分与极大函数}
\label{subsec:ch5-vector-valued-singular-maximal}

向量值推广在调和分析中应用广泛, 因为许多算子可看作向量值奇异积分. 例如 Hardy--Littlewood 极大函数可作如下处理. 取 $A=\mathbb R$, $B=L^\infty(\mathbb R_+)$, 定义 $K(x)=\{\varphi_t(x)\}_{t>0}$, 则相应卷积算子满足 $\|Tf(x)\|_B=M_\varphi f(x)$. 若 $\varphi$ 是满足径向递减控制的光滑近似恒等核, 则 $K$ 满足 \eqref{eq:ch5-vector-hormander-y}, 因而定理~\ref{thm:ch5-vector-valued-calderon-zygmund} 给出 $M_\varphi$ 的 $L^p$ 有界性.

类似地, 把值域取为适当函数空间, 可以处理 Littlewood--Paley 平方函数, 参数族极大算子以及向量值截断奇异积分. 这种方法的一个优点是避免逐个重复标量证明.

\hypertarget{chapter-05-page-124}{}
\subsubsection{强奇异积分}
\label{subsec:ch5-strongly-singular-integrals}

定理~\ref{thm:ch5-calderon-zygmund} 对 Fourier 变换有界的任意缓增分布采用 Hörmander 条件. 若在 Fourier 端施加更强衰减, 可以削弱这一条件. 原型是 $T_{a,b}=K_{a,b}*f$, $0<a<1$, $b>0$, 其中
\[
 \widehat K_{a,b}(\xi)=\varphi(\xi)\frac{e^{i|\xi|^a}}{|\xi|^b},
\]
$\varphi\in C^\infty$ 在原点附近为零, 在某个紧集外恒等于 $1$. 显然 $|\widehat K_{a,b}(\xi)|\le A(1+|\xi|)^{-b}$, 而在原点附近
\[
 |K_{a,b}(x)|\approx\frac B{|x|^{n+\delta}},
 \qquad \delta=\frac{na/2-b}{1-a}.
\]
所以该核一般不满足梯度条件, 也不满足 Hörmander 条件, 但这两个增长条件仍决定 $T_{a,b}$ 在 $L^p$ 上的行为.

\begin{Theorem}
\label{thm:ch5-strongly-singular-model}
算子 $T_{a,b}$ 在 $L^p$ 上有界, 如果
\begin{equation}
\label{eq:ch5-strongly-singular-range}
 \left|\frac12-\frac1p\right|
 <\frac{(b/n)(n/2+\delta)}{b+\delta},
 \qquad \delta=\frac{na/2-b}{1-a}.
\end{equation}
若反向严格不等式成立, 则它无界. 若 \eqref{eq:ch5-strongly-singular-range} 中取等号, 则 $T_{a,b}$ 将 $L^p$ 映入 Lorentz 空间 $L^{p,p'}$.
\end{Theorem}

该有界性由 I. Hirschman 与 S. Wainger 建立, 端点估计归功于 C. Fefferman. 后者是下面一般结果的推论.

\hypertarget{chapter-05-page-125}{}
\begin{Theorem}
\label{thm:ch5-strongly-singular-general}
设 $K$ 是 $\mathbb R^n$ 上具有紧支集的缓增分布, 在原点外与局部可积函数一致. 假设存在 $0<\theta<1$, 使
\[
 |\widehat K(\xi)|\le\frac A{(1+|\xi|)^{n\theta/2}},
 \qquad \xi\in\mathbb R^n,
\]
并且对 $|y|\le1$,
\[
 \int_{|x|>2|y|^{1-\theta}}|K(x-y)-K(x)|\,dx\le B.
\]
则 $Tf=K*f$ 在 $L^p$, $1<p<\infty$, 上有界, 并且满足弱型 $(1,1)$ 不等式.
\end{Theorem}

这一结果后来由 P. Sjölin 推广.

\subsubsection{伪微分算子}
\label{subsec:ch5-pseudodifferential-operators}

伪微分算子推广了微分算子与奇异积分. 形式上定义为
\[
 Tf(x)=\int_{\mathbb R^n}\sigma(x,\xi)\widehat f(\xi)e^{2\pi ix\cdot\xi}\,d\xi,
\]
其中符号 $\sigma$ 是 $\mathbb R^n\times\mathbb R^n$ 上的复值函数. 当 $\sigma(x,\xi)=m(\xi)$ 时, $T$ 是乘子算子; 当 $\sigma(x,\xi)=b(x)$ 时, $T$ 是逐点乘法算子. 一般情形比这两种特例更微妙.

按照符号及其导数的大小分类. 对 $m\in\mathbb Z$, 若
\[
 |D_x^\beta D_\xi^\alpha\sigma(x,\xi)|
 \le C_{\alpha,\beta}(1+|\xi|)^{m-|\alpha|},
\]
则称 $\sigma\in S^m$. 伪微分算子可重写为
\begin{equation}
\label{eq:ch5-pseudodifferential-kernel-form}
 Tf(x)=\int_{\mathbb R^n}K(x,x-y)f(y)\,dy,
\end{equation}
其中 $K$ 由 \eqref{eq:ch4-variable-kernel-inverse-fourier} 给出. 当 $\sigma\in S^0$ 时, $K$ 满足 Hörmander 条件, 且 $T$ 在 $L^2$ 上有界, 因而本章理论给出 $T$ 在 $L^p$, $1<p<\infty$, 上有界.

\hypertarget{chapter-05-page-126}{}
伪微分算子的复合对应于符号微积分. 若 $\sigma_1\in S^{m_1}$, $\sigma_2\in S^{m_2}$, 则相应算子的复合具有 $S^{m_1+m_2}$ 中的符号, 且主项为 $\sigma_1\sigma_2$. 若 $\sigma\in S^m$ 是椭圆的, 即 $|\sigma(x,\xi)|\ge C(1+|\xi|)^m$ 当 $|\xi|$ 足够大时成立, 则存在 $S^{-m}$ 中的参数子, 复合与恒等算子只相差低阶误差. 因而伪微分算子理论可用于椭圆算子求逆.

更一般的符号类 $S_{\rho,\delta}^m$ 由
\[
 |D_x^\beta D_\xi^\alpha\sigma(x,\xi)|
 \le C_{\alpha,\beta}(1+|\xi|)^{m-\rho|\alpha|+\delta|\beta|}
\]
定义. 前述 $S^m$ 即 $S_{1,0}^m$. 当 $m=0$, $0\le\delta<\rho\le1$ 时, 相应核满足标准估计并给出 $L^p$ 有界算子. 关于伪微分算子的进一步结果与参考文献, 见 Stein \MBUnresolvedReference{citation}{[17]}, Journé \MBUnresolvedReference{citation}{[8]} 以及 Coifman 与 Meyer \MBUnresolvedReference{citation}{[2]}.
